%% Figure block, shared by ../paper/main.tex and ../submission/*.tex so the caption has
%% exactly one source.
\begin{figure}[tbp]
\centering
\figorbox{fig1_headline.pdf}
\caption{\textbf{Dependence-blind confidence intervals fail by a computable amount, and
only estimating the long-run score covariance repairs them.}
(a) Coverage of nominal $95\%$ intervals on the homogeneous autoregressive design as the
dependence strength varies, against the parameter-free prediction of
\Cref{prop:coverage}; error bars are $\pm1.96$ Monte Carlo standard errors.
(b) Coverage of every method on that design at $\phi=\psi=0.7$ ($\kappa=\KappaLoArhom$), one row
per method; the dotted line is the nominal level and the number at the right is the interval width
relative to the infeasible oracle interval. Read it in three groups. The five rows at the top all
account for the autocovariances of the score --- the infeasible oracle, the offline reference, and
the three streaming estimators that use our long-run covariance --- and they are the rows that
reach nominal, at width $\approx1$. The next two use the instantaneous covariance the
independent-data theory prescribes; they are narrower by the factor $\sqrt\kappa=\WidthRatioArhom$
and they miss, by the amount \Cref{prop:coverage} predicts. The bottom four are first-order methods
whose \emph{point} estimates have not converged on this ill-conditioned design (their $\sqrt N$
error is $\RmseAdagradArhet$ to $\RmseAsgdArhom$ times the efficient value), so their coverage is
dominated by that and not by which covariance they use; \Cref{tab:wellcond} is where they become
informative about inference.}
\label{fig:coverage}
\end{figure}
